%-----------------------------------------------------------------------------------------------
\makeatletter
\immediate\write18{datelog > \jobname.info} % site script for $(date '+%Y-%m-%d %Hh%Mm%Ss')
\makeatother
%-----------------------------------------------------------------------------------------------
%-----------------------------------------------------------------------------------------------
\usetheme{Copenhagen}
\usepackage{beamercolorthemeCNThermSci2}
\usefonttheme{serif}
%-----------------------------------------------------------------------------------------------

%-----------------------------------------------------------------------------------------------
%-----------------------------------------------------------------------------------------------
\usetheme{Copenhagen}
\usepackage{beamercolorthemeUTF2}
\usefonttheme{serif}
%-----------------------------------------------------------------------------------------------
\usepackage[utf8]{inputenc}
\usepackage[greek,french,english,brazil]{babel} % last becomes the active one
\usepackage{pslatex}
\usepackage{amssymb,amsmath}
\usepackage{soul}
\usepackage[squaren,Gray,cdot]{SIunits}
\usepackage[nice]{nicefrac}
\usepackage{tikz}
\usepackage{amscd}
\usepackage{stmaryrd}
\usepackage{scalerel}
\usepackage{xspace}
%-----------------------------------------------------------------------------------------------


%-----------------------------------------------------------------------------------------------
%-----------------------------------------------------------------------------------------------
% Mathematical
%-----------------------------------------------------------------------------------------------
\newcommand{\vet}[1]{\underline{{#1}}}
\newcommand{\mat}[1]{\underline{\underline{{#1}}}}
\newcommand{\cub}[1]{\underline{\underline{\underline{{#1}}}}}
\newcommand{\eqdef}{{\ensuremath\stackrel{\text{\tiny def}}{=}}}
%-----------------------------------------------------------------------------------------------
% Linguistic
%-----------------------------------------------------------------------------------------------
\newcommand{\GRtxt}[1]{\begin{otherlanguage}{greek}{{#1}}\end{otherlanguage}}
\newcommand{\FRtxt}[1]{\begin{otherlanguage}{french}{{#1}}\end{otherlanguage}}
%-----------------------------------------------------------------------------------------------
% Presentation
%-----------------------------------------------------------------------------------------------
\newcommand{\BkgImgH}[1]{% Places an image centered on the slide background filling the height
    \usebackgroundtemplate{\parbox{\paperwidth}{%
        \vspace*{1sp}\centering\includegraphics[height=\paperheight]{{#1}}
}}}
\newcommand{\BkgImgW}[1]{% Places an image centered on the slide background filling the width
    \usebackgroundtemplate{\parbox{\paperwidth}{%
        \vspace*{1sp}\centering\includegraphics[width=\paperwidth]{{#1}}
}}}
\newcommand{\ArtEndH}[3]{% Transitions to plain image (last) slide: #1:prefix #2,#3:extensions
    \BkgImgH{root/../art/#1.#2}
    \frame<handout:0>[plain]{%
        \transdissolve\vspace*{72mm}\color{white}\scriptsize\bf\input{root/../art/#1.#3}}
    \usebackgroundtemplate{\mbox{~}}
}
\newcommand{\ArtEndW}[3]{% Transitions to plain image (last) slide: #1:prefix #2,#3:extensions
    \BkgImgW{root/../art/#1.#2}
    \frame<handout:0>[plain]{%
        \transdissolve\vspace*{72mm}\color{white}\scriptsize\bf\input{root/../art/#1.#3}}
    \usebackgroundtemplate{\mbox{~}}
}
\newcommand{\ImgColW}[3]{% Inserts a full-width image in a column
    \includegraphics[width=\columnwidth]{root/../art/#1.#2}\\[-0.5\baselineskip]
    \parbox{\columnwidth}{\tiny\hfill\scalebox{0.85}{\input{root/../art/#1.#3}}}
}
\newcommand{\txtpic}[1]{%
    \fcolorbox{lightgray}{white!90!black}{{#1}} 
}
%-----------------------------------------------------------------------------------------------


%-----------------------------------------------------------------------------------------------
\newcommand{\dxyz}[3]{\left(\frac{\partial {#1}}{\partial {#2}}\right)_{#3}} % Display (∂x/∂y)_z
\newcommand{\lxyz}[3]{(\partial {#1}/\partial {#2})_{#3}}                    % Line (∂x/∂y)_z
%-----------------------------------------------------------------------------------------------
\title{A.07.01 -- Relações de Propriedades Termodinâmicas}
\subtitle{Funções Características e Variáveis Reduzidas}
\author{Prof.~C.~Naaktgeboren, PhD}
\date{{\scriptsize\tt%
    \includegraphics[height=6.0mm]{cc/by-nc-nd-88x31.pdf}\\[\smallskipamount]
    https://github.com/CNThermSci/ApplThermSci\\
    Compiled on \input{\jobname.info}
}}
%-----------------------------------------------------------------------------------------------
\begin{document}
%-----------------------------------------------------------------------------------------------
\logo{%
    \parbox{158mm}{% There's a 1mm gap on each side of the 160mm x 90mm slide logo line
        \mode<beamer>{
            \includegraphics[height=6.0mm]{root/00-res/UTFPR/UTFPR-logo-D.pdf}\hfill%
            \includegraphics[height=9.0mm]{root/00-res/logo/CNThermSci-logo-A.pdf}%
        }
        \mode<handout>{
            \includegraphics[height=6.0mm]{root/00-res/UTFPR/UTFPR-logo-W.pdf}\hfill%
            \includegraphics[height=9.0mm]{root/00-res/logo/CNThermSci-logo-W.pdf}%
        }
    }
} % The (delineated, alpha), or washed-out logos
%-----------------------------------------------------------------------------------------------
\frame{\titlepage}
%-----------------------------------------------------------------------------------------------

%-----------------------------------------------------------------------------------------------
\frame{\tableofcontents}
%-----------------------------------------------------------------------------------------------

%-----------------------------------------------------------------------------------------------
\section{Funções Características}
%-----------------------------------------------------------------------------------------------

%-----------------------------------------------------------------------------------------------
\subsection{Possíveis Funções}
%-----------------------------------------------------------------------------------------------

    % !j 96 -i8
    %-------------------------------------------------------------------------------------------
    \begin{frame}{Funções Características}
        \begin{itemize}
            \item Cada \alert{função de Gibbs} estabelece a relação entre três diferenciais;
            \item As tríades: \alert{($du$, $ds$, $dv$)}; \alert{($dh$, $ds$, $dP$)};
                \alert{($da$, $dT$, $dv$)} e \alert{($dg$, $dT$, $dP$)}.
            \item Se \alert{qualquer uma} das funções for conhecida para uma \alert{substânica
                pura}, i.e.,
            \item Se \alert{$f(u, s, v) = 0$}, ou \alert{$f(h, s, P) = 0$}, ou \alert{$f(a, T,
                v) = 0$}, ou \alert{$f(g, T, P) = 0$} o for;
            \item então \alert{todas} as propriedades termodinâmicas\footnote{De engenharia.} da
                substância podem ser obtidas!
            \item Tais funções são conhecidas por \alert{funções características} da substância.
            \item Equivalem a: \alert{$u\!:\!u(s, v)$}, \alert{$h\!:\!h(s, P)$},
                \alert{$a\!:\!a(T, v)$}, ou a \alert{$g\!:\!g(T, P)$}.
        \end{itemize}
    \end{frame}
    %-------------------------------------------------------------------------------------------

    % !j 96 -i8
    %-------------------------------------------------------------------------------------------
    \begin{frame}{Funções Características}
        \alert{Exemplo}: Suponha \alert{$u\!:\!u(s, v)$} seja conhecida para uma substância,
        então:
        %
        \begin{alignat*}{3}
            \uncover<1->{
                \alert{P(s, v)}     &\alert{=      -\dxyz uvs},             &   \qquad  & \mbox{e}\qquad        &   \alert{T(s, v)}     &\alert{=  \dxyz usv},      \qquad\mbox{e}    \\
            }
            \uncover<2->{
                \alert{h(s, v)}     &\alert{=      u(s, v) + P(s, v)v},     &   \qquad  & \rightharpoondown     &   \alert{h(s, v)}     &\alert{=  u - \dxyz uvsv}, \qquad\mbox{e}    \\
            }
            \uncover<3->{
                \alert{a(s, v)}     &\alert{=      u(s, v) - T(s, v)s},     &   \qquad  & \rightharpoondown     &   \alert{a(s, v)}     &\alert{=  u - \dxyz usvs}, \qquad\mbox{e}    \\
            }
            \uncover<4->{
                \alert{g(s, v)}     &\alert{=      h(s, v) - T(s, v)s},     &   \qquad  & \rightharpoondown     &   \alert{g(s, v)}     &\alert{=  u - \dxyz uvsv - \dxyz usvs}.
            }
        \end{alignat*}
    \end{frame}
    %-------------------------------------------------------------------------------------------

%-----------------------------------------------------------------------------------------------
\subsection{Funções Adotadas em Modelagem}
%-----------------------------------------------------------------------------------------------

    % !j 96 -i8
    %-------------------------------------------------------------------------------------------
    \begin{frame}{Modelagem Termodinâmica de Substâncias}

        \begin{itemize}
            \item<1-> Substâncias puras tem suas propriedades termodinâmicas \alert{modeladas}
                via energia de \alert{Helmholtz}, \alert{$a\!:\!a(T, v)$}.
                \vspace*\medskipamount
            \item<2-> Aparentemente, a escolha de \alert{$a\!:\!a(T, v)$} se deve à entropia
                \emph{não} ser parâmetro de entrada, como em \alert{$u\!:\!u(s, v)$} ou
                \alert{$h\!:\!h(s, P)$};
                \vspace*\medskipamount
            \item<3-> Ainda, para \alert{$g\!:\!g(T, P)$}, o par \alert{$(T, P)$} nem sempre é
                \alert{independente} em \alert{substâncias puras}: considere as mudanças de
                fase.
                \vspace*\medskipamount
            \item<4-> Por outro lado, a formulação \alert{$a\!:\!a(T, v)$} baseia-se em duas das
                variáveis \alert{$PvT$}, de \alert{fácil medição}, que \alert{\emph{sempre}} são
                independentes.
        \end{itemize}

    \end{frame}
    %-------------------------------------------------------------------------------------------

    % !j 96 -i8
    %-------------------------------------------------------------------------------------------
    \begin{frame}{Modelagem via energia específica de Helmholtz}
        Se \alert{$a\!:\!a(T, v)$} é uma função conhecida para uma substância, então:
        %
        \begin{alignat*}{3}
            \uncover<1->{
                \alert{P(T, v)}     &\alert{= -\dxyz avT}\quad\mbox{(EoS)!} &   \qquad  & \mbox{e}\qquad        &   \alert{s(T, v)}     &\alert{= -\dxyz aTv},      \qquad\mbox{e}    \\
            }
            \uncover<2->{
                \alert{u(T, v)}     &\alert{=      a(T, v) + Ts(T, v)},     &   \qquad  & \rightharpoondown     &   \alert{u(T, v)}     &\alert{=  a - T\dxyz aTv}, \qquad\mbox{e}    \\
            }
            \uncover<3->{
                \alert{h(T, v)}     &\alert{=      u(T, v) + vP(T, v)},     &   \qquad  & \rightharpoondown     &   \alert{h(T, v)}     &\alert{=  a - T\dxyz aTv - v\dxyz avT}, \qquad\mbox{e}    \\
            }
            \uncover<4->{
                \alert{g(T, v)}     &\alert{=      h(T, v) - Ts(T, v)},     &   \qquad  & \rightharpoondown     &   \alert{g(T, v)}     &\alert{=  a - v\dxyz avT}.
            }
        \end{alignat*}
    \end{frame}
    %-------------------------------------------------------------------------------------------

    % !j 96 -i8
    %-------------------------------------------------------------------------------------------
    \begin{frame}{Funções Características vs. Equações de Estado (EoS)}
        \begin{itemize}
            \item<1-> A equação de estado, \alert{$P\!:\!P(T, v)$} vem de \alert{$-\lxyz avT$};
                \vspace*\bigskipamount
            \item<2-> $\rightharpoondown$ a \alert{função característica} contém \alert{mais
                informações} do que a \alert{equação de estado}!
                \vspace*\bigskipamount
            \item<3-> A entropia, \alert{$s\!:\!s(T, v)$} está `imbutida' em \alert{$a\!:\!a(T,
                v)$}, via \alert{$-\lxyz aTv$};
                \vspace*\bigskipamount
            \item<4-> $\rightharpoondown$ a \alert{função característica} reúne o
                \alert{comportamento $P-v-T$} e o \alert{entrópico/calórico}!
        \end{itemize}
    \end{frame}
    %-------------------------------------------------------------------------------------------

%-----------------------------------------------------------------------------------------------
\section{Variáveis Reduzidas}
%-----------------------------------------------------------------------------------------------

    %------------------------------------------------------------------------------------------
    \begin{frame}[allowframebreaks]{Variáveis Reduzidas}
        Em \alert{ciências térmicas} é muito comum trabalhar com quantidades
        \alert{adimensionais}. \vspace*\medskipamount

        Em \alert{termodinâmica}, tais quantidades são também conhecidas como \alert{reduzidas}.
        \vspace*\medskipamount

        Tais são frequentemente obtidas \alert{dividindo} quantidades pela sua \alert{escala}:
        %
        \begin{alignat*}{5}
            \alert{P^\star} &\alert{\equiv P_{cr}},  & \qquad & \qquad & \alert{T^\star} &\alert{\equiv T_{cr}},          & \qquad & \qquad & \alert{v^\star} &\alert{\equiv v_{cr}}, \\
            \alert{s^\star} &\alert{\equiv RZ_{cr}}, & \qquad & \qquad & \alert{u^\star} &\alert{\equiv T^\star s^\star}, & \qquad & \qquad & \alert{a^\star} &\alert{\equiv u^\star}.
        \end{alignat*}

        Visto que \alert{$s^\star \equiv RZ_{cr}$}, tem-se \alert{$T^\star s^\star = P^\star
        v^\star$}, e assim \alert{$u^\star = h^\star$} e \alert{$g^\star = a^\star$}.

        \pagebreak
        Sejam as novas variáveis reduzidas:
        %
        \begin{alignat}{5}
            \alert{P_r} &\alert{\equiv P/P^\star} & \qquad & \qquad & \alert{T_r} &\alert{\equiv T/T^\star} & \qquad & \qquad & \alert{v_r} &\alert{\equiv v/v^\star}, \\
            \alert{s_r} &\alert{\equiv s/s^\star} & \qquad & \qquad & \alert{u_r} &\alert{\equiv u/u^\star} & \qquad & \qquad & \alert{a_r} &\alert{\equiv a/a^\star}. 
        \end{alignat}

        Derivadas também podem ser adimensionalizadas:
        %
        \begin{align}
            \partial P &= P^\star\partial P_r,\qquad\mbox{and}\qquad \partial T = T^\star\partial T_r\qquad\rightharpoondown \\
            \dxyz PTv &= \frac{P^\star}{T^\star}\dxyz{P_r}{T_r}{v_r} = \frac{P_{cr}}{T_{cr}}\dxyz{P_r}{T_r}{v_r},
            \label{eq:PTV.reduced}
        \end{align}
        %
        \noindent since a process with constant $v$ is also a process with constant $v_r$.
    \end{frame}
    %------------------------------------------------------------------------------------------

    %------------------------------------------------------------------------------------------
    \begin{frame}{Tópicos de Leitura}
        \begin{thebibliography}{Naaktgeboren, C., 2018}
            %-----
            \bibitem[Naaktgeboren, C., 2018]{2018-NaaktgeborenC-UTFPR}
                Naaktgeboren, C.
                \newblock{%
                    {\em Thermodynamic Properties Relations (Handout)}.
                    \alert{Seções 7 e 8.}
                }
                \newblock{\footnotesize Disponibilizado no AVA.}
        \end{thebibliography}
    \end{frame}
    %------------------------------------------------------------------------------------------

    % Finishes with stunning image, with credit
    \ArtEndWSE{pexels-sebastian-voortman-189349}{jpg}{txt}

%-----------------------------------------------------------------------------------------------
\end{document}
%-----------------------------------------------------------------------------------------------

